Quan sát lớp học là một nguồn dữ liệu quan trọng để hiểu cách học sinh tham gia vào hoạt động học tập, cách các hành vi diễn ra trong giờ học và cách môi trường lớp học thay đổi theo thời gian. Trong các nghiên cứu giáo dục, quan sát trực tiếp hoặc quan sát qua video thường được dùng để ghi nhận những hành vi như đọc, viết, trao đổi, cúi đầu, nghiêng người hoặc sử dụng thiết bị. Những ghi nhận này giúp giáo viên và nhà nghiên cứu có thêm bằng chứng khi phân tích động lực lớp học, thay vì chỉ dựa vào điểm số hoặc cảm nhận chủ quan.

Tuy nhiên, quan sát lớp học thủ công đòi hỏi nhiều công sức. Người quan sát cần được huấn luyện, cần thống nhất tiêu chí mã hóa và phải xem lại video trong thời gian dài để ghi nhãn hành vi \,\cite{pianta2009observation,volpe2025classroom}. Khi lớp học có nhiều học sinh xuất hiện đồng thời, nhiều hành vi chồng lấn, che khuất hoặc thay đổi nhanh, việc duy trì độ nhất quán giữa các lần ghi nhãn càng khó hơn. Do đó, các phương pháp thủ công tuy có giá trị về mặt sư phạm nhưng khó mở rộng khi cần phân tích nhiều video hoặc cần phản hồi trong thời gian ngắn.

Sự phát triển của thị giác máy tính mở ra khả năng tự động hóa một phần quy trình quan sát. Nếu hệ thống có thể nhận diện các hành vi quan sát được trong video, kết quả này có thể được dùng để tạo thống kê, tìm các đoạn đáng chú ý và hỗ trợ người dùng xem lại bằng chứng nhanh hơn. Tuy vậy, đầu ra của bộ phát hiện thị giác chỉ mô tả những gì mô hình nhìn thấy trong khung hình; một nhãn như \textit{Read}, \textit{Write}, \textit{Phone}, \textit{Bow}, \textit{Lean} hoặc \textit{Talk} không tự động phản ánh ý định, mức độ tập trung hoặc chất lượng học tập của học sinh \,\cite{fredricks2004engagement}. Vì vậy, một hệ thống phân tích lớp học cần tách rõ dữ kiện quan sát được với các kết luận giáo dục cần bối cảnh và chuyên môn của con người.

Bên cạnh câu hỏi ``hành vi nào xuất hiện'', phân tích lớp học còn cần trả lời câu hỏi ``các hành vi liên hệ với nhau như thế nào theo thời gian''. Ví dụ, một số tín hiệu hành vi có thể thường xuất hiện trước hoặc sau một tín hiệu khác trong cùng đoạn video. Những liên kết này có thể gợi ý các mẫu diễn biến đáng kiểm tra, nhưng không nên được hiểu ngay là quan hệ nhân quả chắc chắn. Trong phạm vi đồ án, các liên kết được phát hiện được diễn giải như quan hệ dự báo hoặc tương quan theo thời gian trên dữ liệu quan sát thụ động.

Từ các yêu cầu trên, đồ án lựa chọn bài toán xây dựng một khung phân tích video lớp học có khả năng đi từ bằng chứng thị giác đến diễn giải thận trọng. Khung này cần nhận diện hành vi trong video, theo dõi và gom các phát hiện thành chuỗi thời gian, khám phá các liên kết dự báo giữa hành vi và tạo báo cáo quan sát có thể truy vết. Mục đích không phải thay thế giáo viên hoặc nhà nghiên cứu, mà là cung cấp một lớp bằng chứng trung gian giúp họ kiểm tra nhanh hơn, có hệ thống hơn và hạn chế diễn giải vượt quá dữ liệu.

Hệ thống được đề xuất trong đồ án có tên CausClass. Tên gọi này nhấn mạnh hai ý: hệ thống làm việc với dữ liệu lớp học, và hệ thống tìm các liên kết có cấu trúc giữa những tín hiệu hành vi. Tuy nhiên, đồ án không khẳng định các cạnh đồ thị là quan hệ nhân quả theo nghĩa can thiệp. CausClass được thiết kế như một công cụ hỗ trợ quan sát dựa trên bằng chứng, trong đó mọi kết quả cuối cùng đều phải có khả năng quay lại video, sự kiện phát hiện, chuỗi thời gian hoặc đồ thị đã kiểm chứng để đối chiếu.